\section*{Capítulo 8: Las ondas gravitatorias: curvatura en movimiento}
\begin{description}[style=multiline,leftmargin=2em,labelwidth=1.5em]

% DOCX paragraph 2086
\item[1.]\hypertarget{note:08-01}{}Una ecuación de ondas para un campo escalar \(\phi\) se puede escribir simplemente como \(\ \square\ \phi\  = \ 0\), donde \(\square\) es el \emph{operador diferencial de d'Alembert\index{d'Alembert, operador de}}, o simplemente dalembertiano, una representación abreviada de una serie de derivadas parciales de segundo orden. Por ejemplo, en un sistema de coordenadas cartesianas de tres dimensiones espaciales, \(\square\  = \partial_{t}^{2} - v^{2}(\partial_{x}^{2} + \partial_{y}^{2} + \partial_{z}^{2})\), donde \(v\) es la velocidad de propagación de la onda.

% DOCX paragraph 2087
\item[2.]\hypertarget{note:08-02}{}Como se ha mencionado en el texto, la descripción del movimiento de la colchoneta es válida a primer orden en la perturbación. Si quisiéramos incluir también efectos no lineales, es decir, contribuciones a la dinámica de la colchoneta de orden superior al primero, esta también se movería en la dirección de propagación de la onda, ya que la onda es parcialmente longitudinal y transfiere cierta cantidad de momento lineal. En la práctica, también nos alejaríamos ligeramente de la persona que saltó al agua.

% DOCX paragraph 2088
\item[3.]\hypertarget{note:08-03}{}Además de estas dos polarizaciones lineales, existen dos \emph{circular\index{polarización!circular}es}, una para cada sentido de giro, en las que la deformación producida por la onda gravitatoria gira alrededor de la dirección de propagación de la onda. Si volvemos a la figura \ref{fig:08-02}, en este caso el Hombre de Vitruvio\index{Leonardo da Vinci!Hombre de Vitruvio} se estiraría y comprimiría en una dirección que va girando con el tiempo mientras se propaga la onda gravitatoria.

% DOCX paragraph 2089
\item[4.]\hypertarget{note:08-04}{}En los primeros instantes de su vida, el universo era extremadamente caliente y denso, de modo que la materia y la radiación estaban «acopladas» y sus temperaturas evolucionaban de la misma manera. Cuando el universo se expandió y se enfrió lo suficiente, la materia y la radiación se «desacoplaron» y, en términos de evolución térmica, cada una «siguió su propio camino». El momento del desacoplamiento fue muy especial, pues los fotones quedaron por fin «libres» para propagarse sin dispersarse constantemente al encontrarse con las partículas de su entorno. Estos fotones, producidos tan solo 380\,000 años después del Big Bang, siguen siendo visibles en el llamado \emph{fondo cósmico de microondas\index{fondo cósmico de microondas}}, un baño de radiación muy uniforme. Si nuestros ojos fueran sensibles a las microondas, el cielo nocturno no parecería oscuro, sino tenuemente iluminado por esa radiación primordial.

% DOCX paragraph 2090
\item[5.]\hypertarget{note:08-05}{}Aunque es fácil definir un horizonte de sucesos en un espaciotiempo estático o estacionario, calcularlo es mucho más difícil en un espaciotiempo dinámico, como el de dos agujeros negros poco antes de fusionarse o el de una estrella antes de colapsar en un agujero negro. En estos casos resulta mucho más conveniente adoptar el concepto de \emph{«horizonte aparente»} para definir la superficie bidimensional dentro de la cual la luz queda atrapada en un instante dado. En un espaciotiempo estático o estacionario, los dos tipos de horizonte coinciden.

% DOCX paragraph 2091
\item[6.]\hypertarget{note:08-06}{}El número de Love no tiene nada que ver con los sentimientos: lleva el nombre del físico británico Augustus\index{Love!Augustus} Love (1863--1940), quien lo definió por primera vez en su teoría de la elasticidad.

% DOCX paragraph 2092
\item[7.]\hypertarget{note:08-07}{}Por supuesto, si la masa del objeto supera unas 3 masas solares, no será una estrella de neutrones, sino un agujero negro. Sin embargo, hay observaciones en las que la distinción no resulta tan clara. Por ejemplo, en GW190814\index{GW190814}, el objeto menos masivo tiene unas 2,6 masas solares. Por tanto, no está claro si las observaciones corresponden al agujero negro menos masivo jamás observado o a la estrella de neutrones más masiva detectada hasta ahora; véase The LIGO Collaboration and Virgo Collaboration, en \emph{Astrophysical Journal Letters}, 896, L44 (2020).

% DOCX paragraph 2093
\item[8.]\hypertarget{note:08-08}{}Buena parte de lo que hemos visto para sistemas binarios de agujeros negros o estrellas de neutrones también se aplica a un sistema binario mixto, compuesto por un agujero negro y una estrella de neutrones. Sin embargo, hay dos diferencias importantes. La primera está relacionada con la razón de masas entre ambos objetos, que puede ser de 10 a 1 o incluso mayor, dadas sus masas características: se espera que el agujero negro tenga entre 15 y 25 masas solares y la estrella de neutrones entre 1,3 y 2,3. La segunda diferencia es que la fusión de un sistema mixto no puede dar lugar a una estrella hipermasiva. En este caso, la estrella de neutrones es destruida por las fuerzas de marea y forma un toro de acreción, o incluso es «absorbida entera» si la diferencia de masas es muy grande.

% DOCX paragraph 2094
\item[9.]\hypertarget{note:08-09}{}Véase E. Nakar, en \emph{Physics Reports}, 442, 166 (2007), para una revisión de la bibliografía sobre este tema.

% DOCX paragraph 2095
\item[10.]\hypertarget{note:08-10}{}Véase S. E. Woosley y J. S. Bloom, en \emph{Annual Review of Astronomy and Astrophysics}, 44, 507 (2006), para una revisión de la bibliografía pertinente.

% DOCX paragraph 2096
\item[11.]\hypertarget{note:08-11}{}Los detectores de ondas gravitatorias son, en general, «\emph{omnidireccionales}», es decir, sensibles en todas las direcciones, pero no con la misma sensibilidad. De hecho, hay direcciones en las que la sensibilidad es máxima y otras en las que es muy pequeña, los «puntos ciegos» del detector. Por tanto, como ocurre con una antena de radio, se puede definir un mapa de sensibilidad de un detector interferométrico de ondas gravitatorias. En el caso de GW170817\index{GW170817}, el hecho de que Virgo no midiera la señal resultó, aun así, muy valioso. Conocer su punto ciego en aquel momento permitió restringir la región de búsqueda de la señal electromagnética y, por tanto, descubrir GRB 170817A\index{GRB 170817A}.

% DOCX paragraph 2097
\item[12.]\hypertarget{note:08-12}{}Por lo general, no se espera una contrapartida electromagnética\index{contrapartida electromagnética} intensa de la fusión de sistemas binarios de agujeros negros de masa estelar. La razón es que estos sistemas representan la etapa final de un sistema binario que hace mucho perdió todo rastro de materia, ya fuera mediante los procesos de colapso que produjeron los agujeros negros o mediante las explosiones de supernova y las ondas de choque\index{choque} asociadas, que dispersaron la materia por el espacio interestelar.

% DOCX paragraph 2098
\item[13.]\hypertarget{note:08-13}{}L. Rezzolla, B. Giacomazzo, L. Baiotti, J. Granot, K. Koveliotou, M.-A. Aloy, en \emph{The Astrophysical Journal Letters}, 732, L6 (2011).

% DOCX paragraph 2099
\item[14.]\hypertarget{note:08-14}{}L.-X. Li y B. Paczyński, en \emph{The Astrophysical Journal Letters}, 507, L59 (1998).

% DOCX paragraph 2100
\item[15.]\hypertarget{note:08-15}{}En astronomía se llama \emph{nova\index{nova}} a una explosión termonuclear en la superficie de una enana blanca, provocada por la acumulación de hidrógeno debido a la acreción de materia de una estrella compañera en un sistema binario. La explosión puede ser visible durante varios días y alcanza una luminosidad de entre 10\,000 y 100\,000 veces la del Sol. Estas explosiones liberan cantidades diminutas de energía en comparación con los brotes de rayos gamma.

% DOCX paragraph 2101
\item[16.]\hypertarget{note:08-16}{}J. M. Lattimer y D. N. Schramm, en \emph{The Astrophysical Journal}, 210, 549 (1976).

\end{description}
